% =====================================================================
%  abstract.tex —— 摘要
% =====================================================================

\chapter*{摘要}
\addcontentsline{toc}{chapter}{摘要}
\markboth{摘要}{摘要}

本文档系统记述 \textbf{Huntian} —— 一套面向 GaN 射频高电子迁移率晶体管
（high electron mobility transistor, HEMT）的轻量化、模块化工艺计算机辅助设计
（technology computer-aided design, TCAD）内核 —— 的计算原理、离散格式与验证证据。

Huntian 的架构主张是\emph{模块即插件}：每一个物理模型、每一个数值服务都以独立
插件的形式挂在注册中心上，彼此之间不通过代码调用而只通过命名数据通道
（artifact）交换数据。因此，本文档采用与之同构的组织方式：\textbf{一个插件对应一个
模块章}，各章自洽地给出该模块的控制方程、离散化、算法、参数依据与验证基准，
而架构层与验证方法学集中在前置的总论部分。

每章遵循固定的六段式骨架：模块契约、控制方程与物理假设、离散化与数值算法、
可标定参数、验证与基准、误差来源与适用边界。这种一致性并非形式主义 ——
它使读者可以在任意两章之间横向比较\emph{同一类证据}，也使文档本身可以被机械化
校验（章首的模块标识与代码中的 \texttt{MANIFEST} 逐字段比对）。

本文档对验证证据采取严格的二分表述：凡进入"实测"列的数字，均可在仓库的测试
套件或基准记录中回查；凡未经验证者，一律在"适用边界"一节明确列出，不作粉饰。
截至本文档当前版本，仓库测试套件共 \textbf{426} 个用例，全部通过。

\vspace{1em}
\noindent\textbf{关键词：} GaN HEMT；TCAD；自洽 Schrödinger--Poisson 方程；
二维电子气；漂移扩散；Scharfetter--Gummel 离散；Newton 法；精确 Jacobian；
插件架构；可复现性

\cleardoublepage
